\begin{tcolorbox}
\small
Thanks for providing feedback on our AI4Science benchmark called ScienceAgentBench.
We are developing an AI agent to assist you! Given a task instruction and a dataset, the agent will help you write a computer program to fulfill the task you have in mind.
To develop and evaluate such an AI agent, we have collected a benchmark by adapting some tasks from peer-reviewed publications with open-source codes. Each data sample in our benchmark consists of the following main components:

\textbf{Task Instruction}: Describes (1) the goal of a task or a scientific hypothesis and (2) output requirements.

\textbf{Dataset Information}: Contains (1) the dataset directory structure and (2) helpful metadata or a few examples from the dataset.

\textbf{Annotated Program}: The reference solution adapted from each publication's open-source code.

\textbf{Evaluation Script}: The code to evaluate AI agents' performance by comparing the execution results of its generated programs with those of the annotated programs.

To ensure that each task is formulated and described correctly and professionally, we would like you to give us a hand by reviewing our collected data samples. In addition, we are also seeking some additional information from you as a domain expert, including writing down some task-related domain knowledge and revising a rubric to score the generated programs.

Please follow the guidelines below to review each task. First, please enter the Task ID you are reviewing: [task\_id]

\textbf{Guidelines for Data Reviewing}

First, a bit more background: once you give the AI agent a task instruction, it will try to automatically complete everything without seeking additional help from you. This is similar to the scenario where you give the task to a junior student in your lab/class who will complete it as an assignment.

For each task, please first spend a few minutes reading the given task information (instruction, dataset information, and source GitHub repository) and our annotated program to have a rough understanding of the task and relevant concepts. Then, please comment on the following two parts. Note that you may iteratively revise your answer to each question to help us improve the task instructions and programs.

\textbf{1. Program}

Is the program a valid solution (not necessarily the best solution) to the given task instruction? Here is an example: [google\_doc\_link] \footnote{The example for program and instruction validation is provided in Section \ref{app:program_example}.}

If there are only minor issues, please comment on how the program should be modified below. 

However, if you believe there is a major issue (e.g., the program is doing sth irrelevant or more than two lines of code need to be revised in order to make it correct), please let us know the task ID and do NOT fill the rest of the form.

Is the program a valid solution to the given task instruction?

[] Yes

[] Need Modification (comment below)

[] No (report and continue to the next task)

How should the program be modified? Please mention the line numbers that need to be inspected.

[Long Text Answer]

\textbf{2. Task Instruction}

The task instructions were created by non-experts and thus might contain some misused terms,  awkward expressions, or inaccurate descriptions of the task that do not adhere to your domain’s scientific language. Do you see such issues for this task instruction? If so, please revise or rewrite the task instruction for any issues you can find.

Finally, if needed, please help make the task instruction more fluent and natural sounding. 

Please enter your revised task instruction below. If there are no changes, please skip this question and leave the answer text blank. 

[Long Text Answer]
\end{tcolorbox}

\begin{tcolorbox}
\small
\textbf{3. Domain Knowledge}

Suppose the AI agent fails to fulfill the task based solely on the task instruction, perhaps due to lack of some background knowledge, we want to provide some additional information to help it succeed. This is similar to the situation where you give an exam problem to a student in your class and they might not be able to do it just based on the problem description, but you can provide some hints to help them. 

Please write down at most three important pieces of knowledge that are related to the task and the program. 

For example ([google\_doc\_link]): \footnote{The example for domain knowledge annotation is provided in Section \ref{app:knowledge_example}.}

Concepts and details in the task description or program that may need further explanation or extra attention, e.g., a term definition on wikipedia you would send to a new student (without much domain expertise) working on this task, or a common practice for such tasks in your field.

Information about the python packages and/or functions used in the program. For example,  you would copy and paste a snippet of package description/function documentation to help the new student working on the task. 

You may assume the AI agent has a general sense of your domain, like a new graduate student with undergrad-level knowledge but not much about the specific task. Please help it by providing some knowledge to write the program for this task. You can search online for more details about the dataset, packages, and functions used in each task before writing. 

Please try not to “leak” the annotated program directly. You may imagine that you don't have the direct answer but could provide some helpful information to your junior colleague so that they can derive the program. For example:

Instead of copying/describing a few lines in the program, you may copy the documentations of packages/functions used in that program.

Instead of specifying the variables and parameters, you may suggest a range (e.g., 1e-3 to 1e-4 for learning rate).

Instead of saying columns A,B,C are related to the target attribute Y, you may try to find a knowledge snippet describing what is correlated to Y.

However, in some rare cases, there may be a need to provide a minimal “leak” of the annotated program, e.g., the decision boundary of Y is 0.6 instead of 0.5. Still, it would be great if you could think about its necessity before annotating such knowledge.

For each piece of domain knowledge related to this task, please write 1-5 sentences. If you believe the task instruction is self-contained and needs no further explanations, please enter "None".

[Long Text Answer]

\textbf{4. Scoring Rubric}

Once the AI agent generates a program, we need an evaluation method to review the generated program. To do it, we need a task-specific scoring rubric, which assigns partial credits for more comprehensive evaluation of the generated program. Right now we have already got an initial draft of the rubric with five major components: (1) data loading, (2) data processing, (3) modeling, analysis or visualization, (4) output formatting, (5) saving output.

Please review our initial draft of the rubric. Imagine that you will use this rubric to score the programs produced by your junior students. Please modify the rubric items that you think are incorrect, should be described with more/less details, or should be reweighed with higher/lower credits for each component. Please also add any missing but necessary rubric item you would use to assess a program’s correctness, or remove redundant rubric items.

Please enter your revised scoring rubric below. If there are no changes, please skip this question and leave the answer text blank.

[Long Text Answer]
\end{tcolorbox}